% Generated by tools/textbook/export_listings.js from tools/decks/mastery.js and tools/concepts/registry.json: do not edit.
\begin{chapstatus}{Stato di apprendimento di questo capitolo}
\masterychip{now}{DA COMPRENDERE OGGI}\enspace il progetto Verm1; contratto init/update/draw; un colore è un numero; gli array in memoria; .text, .rodata, .data; indirizzo, valore, load, store; perché si salva ra; shift come moltiplicazione; la virgola fissa; le pseudo-istruzioni; l'ordine little-endian.\par\smallskip
\masterychip{use}{DA UTILIZZARE OGGI}\enspace gli strumenti di PRG32, il debugger di PRG32-QT (si usa, non si studia, in questo corso).\par\smallskip
\masterychip{later}{DA COMPRENDERE PIÙ AVANTI}\enspace formati delle istruzioni (Capitolo~\ref{ch:instructions}).
\end{chapstatus}
\begin{chapstatus}{Percorso del capitolo}
Alla fine del capitolo saprete:
\begin{enumerate}[nosep,leftmargin=2em]
\item costruire una cartuccia da un file vuoto: una griglia, un titolo e una testa che si muove
\item mettere costanti e variabili nella sezione giusta, .rodata o .data
\item controllare una velocità con un accumulatore in virgola fissa
\end{enumerate}
Si appoggia su: contratto init/update/draw (Capitolo~\ref{ch:firstcontact}); un colore è un numero (Capitolo~\ref{ch:information}); gli array in memoria (Capitolo~\ref{ch:information}); indirizzo, valore, load, store (Capitolo~\ref{ch:firstcontact}); perché si salva ra (Capitolo~\ref{ch:firstcontact}). Prossimo passo: Capitolo~\ref{ch:boole1}.
\par\smallskip\textbf{Guardando indietro, guardando avanti.} Nella Lezione 02 il passo 2 salvava ra sullo stack prima delle sue chiamate, e lw t1, 0(t0) leggeva una parola all'indirizzo in t0. Oggi verm1\_draw salva ra e verm1\_update no, e ogni variabile di Verm1 si legge con la e lw. La Lezione 18 spiega lo stack e la convenzione di chiamata; la Lezione 19 nomina i modi di indirizzamento. \emph{Richiamo:} Perché verm1\_draw salva ra mentre verm1\_update no?
\end{chapstatus}
